\chapter{Sistemas LIT y la convolución}

\section{Introducción}

Hemos visto que la suma de convolución 
\begin{equation}
	y[n] = \sum_{k=-\infty}^{\infty} x[k] \, h[n-k] = x[n] * h[n] ,
\end{equation}
define la respuesta $y$ de un sistema LIT para una entrada $x$ si conocemos la respuesta al impulso $h$. Por lo tanto, en este capítulo usaremos la convolución como herramienta para estudiar distintas propiedades de los sistemas LIT.

\section{Propiedades de los sistemas LIT}

\subsection{La propiedad conmutativa}

El resultado de la convolución es independiente del ordenamiento de las señales, o en otras palabras, un sistema LIT de respuesta al impulso $h$ excitado con una entrada $x$ genera una salida similar al de un sistema LIT de respuesta al impulso $x$ excitado con una entrada $h$.

Consideremos la operación de convolución:
\begin{equation}
	y[n] = \sum_{k=-\infty}^{\infty} x[k] \, h[n-k] = \sum_{k=-\infty}^{\infty} h[k] \, x[n-k].
\end{equation}

Realizamos un cambio de variable $m = n - k$ y por consecuencia definimos $k = n - m$ para obtener:
\begin{equation}
	y[n] = \sum_{k=-\infty}^{\infty} x[k] \, h[n-k] = \sum_{m=-\infty}^{\infty} x[n-m] \, h[m].
\end{equation}

\subsection{La propiedad distributiva}

La salida de un sistema LIT cuya respuesta al impulso es $h_{a} + h_{b}$ para una entrada $x$ es similar a la suma de las salidas de dos sistemas de respuestas al impulso $h_{a}$ y $h_{b}$ frente a una entrada $x$.

Probaremos la distributividad de la convolución discreta:
\begin{align}
	y[n] &= x[n] * \left( h_{a}[n] + h_{b}[n] \right) , \\
	y[n] &= x[n] * h_{a}[n] + x[n] * h_{b}[n],
\end{align}
planteando:
\begin{equation}
	y[n] = \sum_{k=-\infty}^{\infty} x[k] \, \left( h_{a}[n-k] + h_{b}[n-k] \right).
\end{equation}
Dado que la multiplicación es distributiva respecto a la suma, obtenemos:
\begin{equation}
	y[n] = \sum_{k=-\infty}^{\infty} \left( x[k] \, h_{a}[n-k] + x[k] \, h_{b}[n-k] \right).
\end{equation}
Separando sumatorias:
\begin{equation}
	y[n] = \sum_{k=-\infty}^{\infty} \left( x[k] \, h_{a}[n-k] \right) + \sum_{k=-\infty}^{\infty} \left( x[k] \, h_{b}[n-k] \right) ,
\end{equation}
obtenemos:
\begin{equation}
	y[n] = x[n] * h_a[n] + x[n] * h_{b}[n] .
\end{equation}

\subsection{La propiedad asociativa}

La operación de convolución es asociativa:
\begin{align}
	y[n] &= x[n] * \left( h_a[n] * h_{b}[n] \right) , \\
	y[n] &= \left( x[n] * h_a[n] \right) * h_{b}[n] , \\
	y[n] &= x[n] * h_a[n] * h_{b}[n].
\end{align}

Para probar la propiedad asociativa, planteamos las dos formas de hacer el cálculo. Primero consideremos:
\begin{align}
	y_{1}[n] &= \sum_{p=-\infty}^{\infty} x[p] \, h_{a}[n-p], \\
	y[n] &= \sum_{q=-\infty}^{\infty} y_{1}[q] \, h_{b}[n-q].
\end{align}
Sustituyendo $y_1[q]$:
\begin{equation}
	y[n] = \sum_{q=-\infty}^{\infty} \left( \sum_{p=-\infty}^{\infty} x[p] \, h_{a}[q-p] \right) \, h_{b}[n-q],
\end{equation}
y reordenando:
\begin{equation}
	y[n] = \sum_{q=-\infty}^{\infty} \sum_{p=-\infty}^{\infty} x[p] \, h_{a}[q-p] \, h_{b}[n-q]. \label{eq:asoc1}
\end{equation}

Después consideramos el cálculo alternativo:
\begin{align}
	h[n] &= \sum_{r=-\infty}^{\infty} h_{a}[r] \, h_{b}[n-r], \\
	y[n] &= \sum_{s=-\infty}^{\infty} x[s] \, h[n-s].
\end{align}
Sustituyendo $h[n-s]$:
\begin{equation}
	y[n] = \sum_{s=-\infty}^{\infty} x[s] \, \left( \sum_{r=-\infty}^{\infty} h_{a}[r] \, h_{b}[n - s - r] \right),
\end{equation}
y reordenando:
\begin{equation}
	y[n] = \sum_{s=-\infty}^{\infty} \sum_{r=-\infty}^{\infty} x[s] \, h_{a}[r] \, h_{b}[n - s - r]. \label{eq:asoc2}
\end{equation}

Comparando (\ref{eq:asoc1}) y (\ref{eq:asoc2}), y considerando los cambios de variable $p = s$ y $q - p = r$ (lo que implica $q = r + s$), vemos que:
\begin{equation}
	n - q = n - (r + s) = n - s - r, 
\end{equation}
por lo tanto, ambas expresiones son idénticas, demostrando la propiedad.

\subsection{Preguntas}

\begin{enumerate}
	\item ¿Qué significado tiene la propiedad conmutativa de los sistemas lineales?
	\item ¿Qué significado tiene la propiedad distributiva de los sistemas lineales?
	\item ¿Qué significado tiene la propiedad asociativa de los sistemas lineales?
	\item Dé algún ejemplo concreto de estas tres propiedades.
\end{enumerate}

\subsection{Sistemas LIT con y sin memoria}

Un sistema no tiene memoria si su salida en cualquier momento depende solamente de la entrada de ese momento. Si un sistema no tiene memoria entonces debe ser $h[n]=0$ para todo $n \neq 0$, y por lo tanto:
\begin{equation}
	h[n] = K \, \delta[n] .
\end{equation}
Si un sistema tiene una respuesta al impulso $h[n]$ que es distinta de $0$ para $n \neq 0$, entonces el sistema tiene memoria.

\subsection{Preguntas}

\begin{enumerate}
	\item ¿Conoce algún sistema físico que no tenga memoria?
	\item ¿Conoce algún sistema físico que tenga memoria?
\end{enumerate}

\subsection{Sistemas LIT causales}

Si $y[n_{0}]$ no depende de $x[n]$ para $n > n_{0}$, entonces un sistema LIT es causal. Entonces todos los valores $h[n]$ para $n < 0$ deben ser $0$ y por lo tanto en la suma de convolución podemos cambiar el límite superior:
\begin{equation}
	y[n] = \sum_{k=-\infty}^{n} x[k] \, h[n-k],
\end{equation}
porque $h[n-k] = 0$ cuando $n - k < 0$, es decir, cuando $k > n$.

Además, si la entrada $x[n]$ comienza en $n = 0$ (es decir, $x[n]=0$ para $n<0$), podemos cambiar el límite inferior de la sumatoria:
\begin{equation}
	y[n] = \sum_{k=0}^{n} x[k] \, h[n-k].
\end{equation}
Esta fórmula significa que podemos calcular la salida de un sistema LIT causal con entrada causal de tiempo discreto con una cantidad finita de multiplicaciones y sumas.

\subsection{Preguntas}

Indicar si las siguientes afirmaciones son verdaderas o falsas, en el caso de un sistema lineal invariante en el tiempo:

\begin{enumerate}
	\item Si la entrada es nula para todo tiempo, entonces la salida es nula para todo tiempo.
	\item Si la relación entre entrada y salida es $y[n] = x[n] \, \cos(n)$, entonces el sistema es LIT.
	\item Si la relación entre entrada y salida es $y[n] = \sum_{k=0}^{n} x[k]$, entonces el sistema es LIT.
\end{enumerate}

\subsection{Sistemas LIT estables e inestables}

\subsubsection{Estabilidad BIBO}

El criterio de estabilidad \textbf{BIBO} (Bounded Input Bounded Output o entrada acotada, salida acotada) se puede expresar matemáticamente estableciendo que si la entrada está acotada por una constante $M_x < \infty$:
\begin{equation}
	|x[n]| \leq M_{x},
\end{equation}
entonces la salida debe estar acotada por una constante finita $M_y < \infty$:
\begin{equation}
	|y[n]| \leq M_{y}.
\end{equation}

Sustituyendo la suma de convolución y aplicando la desigualdad triangular (el valor absoluto de una suma es menor o igual a la suma de los valores absolutos), obtenemos:
\begin{equation}
	|y[n]| = \left| \sum_{k=-\infty}^{\infty} x[k] \, h[n-k] \right| \leq \sum_{k=-\infty}^{\infty} |x[k]| \, |h[n-k]|.
\end{equation}
Dado que $|x[k]| \leq M_x$ para todo $k$, podemos extraer la cota de la sumatoria:
\begin{equation}
	|y[n]| \leq M_{x} \sum_{k=-\infty}^{\infty} |h[n-k]|.
\end{equation}
Realizando el cambio de variable $m = n-k$, la sumatoria de $|h[n-k]|$ es idéntica a la sumatoria de $|h[m]|$. Para que la salida esté siempre acotada ($|y[n]| \leq M_{y} < \infty$), se debe cumplir que la suma de los valores absolutos de la respuesta al impulso sea finita:
\begin{equation}
	\sum_{m=-\infty}^{\infty} |h[m]| < \infty.
\end{equation}
A esto se lo conoce como una respuesta al impulso "absolutamente sumable".

\subsection{Respuesta de un sistema LIT en tiempo discreto a una entrada exponencial compleja}

\begin{enumerate}
	\item Vamos a calcular la convolución de una entrada $x[n] = e^{j \, \omega \, n}$ con una respuesta al impulso $h[n]$.
	\item Comenzamos escribiendo la fórmula alternativa de la suma de convolución:
	\begin{equation}
		y[n] = \sum_{k=-\infty}^{\infty} h[k] \, x[n-k].
	\end{equation}
	\item Sustituyendo $x[n - k] =  e^{j \, \omega \, (n-k)}$ obtenemos:
	\begin{equation}
		y[n] = \sum_{k=-\infty}^{\infty} h[k] \, e^{j \, \omega \, (n-k)} = \sum_{k=-\infty}^{\infty} h[k] \, e^{-j \, \omega \, k} \, e^{j \, \omega \, n}. \label{dt.eigenfunc.1}
	\end{equation}
	\item Viendo la ecuación (\ref{dt.eigenfunc.1}) podemos extraer el factor común $e^{j \, \omega \, n}$ de la sumatoria (ya que no depende de $k$) y escribir:
	\begin{equation}
		y[n] = \left(\sum_{k=-\infty}^{\infty} h[k] \, e^{-j \, \omega \, k}\right) \, e^{j \, \omega \, n}. \label{dt.eigenfunc.2}
	\end{equation}
	\item Viendo la ecuación (\ref{dt.eigenfunc.2}), podemos definir $H(e^{j \, \omega})$, la transformada de Fourier en tiempo discreto de $h[n]$:
	\begin{equation}
		H(e^{j \, \omega}) = \sum_{k=-\infty}^{\infty} h[k] \, e^{-j \, \omega \, k}. \label{dt.eigenfunc.3}
	\end{equation}	 
	\item Entonces resulta:
	\begin{equation}
		y[n] = H(e^{j \, \omega}) \, e^{j \, \omega \, n}.
	\end{equation}
	Se dice que las funciones de la forma $e^{j \, \omega \, n}$ son \textit{autofunciones} (o funciones propias) de los sistemas lineales invariantes en el tiempo, siendo $H(e^{j \, \omega})$ su autovalor correspondiente.
\end{enumerate}  

\subsection{Respuestas}

Estas son las respuestas a las preguntas planteadas arriba:

\begin{enumerate}
	\item ¿Qué significado tiene la propiedad conmutativa de los sistemas lineales?
	Un sistema LIT con respuesta al impulso $h[n]$ responde a una entrada $x[n]$ igual que un sistema LIT con respuesta al impulso $x[n]$ responde a una entrada $h[n]$. Además, cuando conectamos en serie dos sistemas LIT, no importa el orden en que se conecten: el conjunto de los dos sistemas va a responder igual ante cualquier entrada.
	
	\item ¿Qué significado tiene la propiedad distributiva de los sistemas lineales?
	Podemos implementar un sistema LIT con respuesta al impulso $h_{1}[n] + h_{2}[n]$ conectando en paralelo dos sistemas LIT cuyas respectivas respuestas al impulso sean $h_{1}[n]$ y $h_{2}[n]$.
	
	\item ¿Qué significado tiene la propiedad asociativa de los sistemas lineales?
	Cuando conectamos en cascada (serie) sistemas LIT, podemos agruparlos para formar nuevos sistemas LIT equivalentes sin alterar la salida final.
	
	\item ¿Conoce algún sistema físico que no tenga memoria?
	Un cable o un divisor resistivo ideal.
	
	\item ¿Conoce algún sistema físico que tenga memoria?
	Un tanque de agua (acumula o "recuerda" el agua que entra) o un capacitor en un circuito eléctrico.
	
	\item ¿La siguiente afirmación es verdadera o falsa? "Si la entrada es nula para todo tiempo, entonces la salida es nula para todo tiempo."
	\textbf{Verdadero}, pues si la entrada es nula, todos los productos en la suma de convolución son nulos.
	
	\item ¿La siguiente afirmación es verdadera o falsa? "Si la relación entre entrada y salida es $y[n] = x[n] \, \cos(n)$, entonces el sistema es LIT."
	\textbf{Falso}, porque el sistema varía en el tiempo. La respuesta a un impulso depende del instante de tiempo en el que este ocurra:
	\begin{equation}
		y[n] = \delta[n - n_{0}] \, \cos(n_{0}).
	\end{equation}
	
	\item ¿La siguiente afirmación es verdadera o falsa? Si la relación entre entrada y salida es $y[n] = \sum_{k=0}^{n} x[k]$, entonces el sistema es LIT.
	\textbf{Verdadero}, pues dicha relación puede escribirse como la convolución de la entrada con un escalón unitario $h[n] = u[n]$:
	\begin{equation}
		y[n] = \sum_{k=-\infty}^{\infty} x[k] \, u[k] \, u[n-k].
	\end{equation}
\end{enumerate}